\section{Introducción}

La geometría fractal permite el estudio de estructuras complejas y 
autosimilares mediante procesos iterativos. Esta investigación aborda 
la transición de estos sistemas desde el plano bidimensional hacia el 
espacio tridimensional mediante los siguientes enfoques:

 \begin{itemize}
    \item \textbf{Triángulo de Sierpinski}: Construcción bidimensional iterativa 
    generada al reducir un triángulo equilátero a su cuarta parte y 
    posicionar tres figuras a escala, formando un patrón continuo de 
    autosimilitud.

    \item \textbf{Tetraedro de Sierpinski y dimensión fractal}: Generalización 
    tridimensional mediante un algoritmo que ensambla cuatro copias de un tetraedro 
    regular con sus aristas reducidas a la mitad. Esta extensión resalta el papel de 
    la dimensión fractal para describir la complejidad espacial y facilita el 
    análisis de la evolución del área hacia el volumen.

    \item \textbf{Construcción Computacional:} Diseño de un pseudocódigo basado 
    en llamadas recursivas que utiliza como parámetros la longitud de la arista, 
    el punto inicial y el número de iteraciones para visualizar matemáticamente 
    la estructura generada.
 \end{itemize}